% Unit 083 terjemahan Bahasa Indonesia dari chapter6.tex baris 2124--2312.
% Sumber: Methods of Algebra, Volume 2, commit
% 9a5803ff77dd3257484cb177f851a73770a59dd3 (CC BY 4.0).
% stable-unit-id: o014.aljabr2.chapter6.tate-cohomology
% Cakupan: Bagian 6.11 lengkap tentang kohomologi Tate; berhenti tepat
% sebelum Bagian 6.12 tentang kohomologi grup profinit pada baris 2313.

% segment-id: o014.li2.chapter6.tate-cohomology.h001
\section{Kohomologi Tate}\label{sec:Tate-coh}
\hypertarget{o014.aljabr2.chapter6.tate-cohomology}{ }

% segment-id: o014.li2.chapter6.tate-cohomology.p001
Untuk grup hingga $G$ dan gelanggang komutatif $\Bbbk$,
Definisi--Proposisi~\ref{def:group-norm} memperkenalkan elemen
$\nu = \sum_{g \in G} g$ dari $\Bbbk[G]^G$. Elemen ini terletak di pusat
$\Bbbk[G]$. Karena itu, untuk setiap modul-$G$ $M$, pemetaan perkalian
$x \mapsto \nu x$ memberikan homomorfisme modul-$G$ $M \to M$, yang tetap
ditulis $\nu$.

% segment-id: o014.li2.chapter6.tate-cohomology.p002
Dari identitas $\nu g = \nu = g \nu$, diperoleh
$\nu M \subset M^G$ dan $\mathfrak{I}M \subset M^{\nu = 0}$, dengan
$\mathfrak{I}$ ideal augmentasi dari $\Bbbk[G]$ dan $M^{\nu=0}$ sebagaimana
dalam Konvensi~\ref{con:G-mod-equalizer}. Dengan demikian, $\nu$ menginduksi
homomorfisme modul-$\Bbbk$
% segment-id: o014.li2.chapter6.tate-cohomology.e001
\begin{equation}\label{eqn:nu-coinv-inv}
	\nu: M_G \to M^G .
\end{equation}

% segment-id: o014.li2.chapter6.tate-cohomology.p003
Karena homologi (atau kohomologi) grup dapat dipandang sebagai versi turunan
dari $M_G$ (atau $M^G$), homomorfisme kanonik $M_G \to M^G$ mengisyaratkan
adanya jembatan antara homologi dan kohomologi. Misalnya, mengapa tidak
meninjau kokernelnya dalam suatu arti turunan\footnote{Dari sudut pandang
topologis: dalam arti homotopi.}? Untuk sementara kita kesampingkan gagasan
ini dan mendefinisikan kohomologi terkait dengan cara klasik.

% segment-id: o014.li2.chapter6.tate-cohomology.d001
\begin{definition}[Kohomologi Tate]\label{def:TaHm}
% segment-id: o014.li2.chapter6.tate-cohomology.x001
	\index{kohomologi Tate (Tate cohomology)}
% segment-id: o014.li2.chapter6.tate-cohomology.x002
	\index[sym1]{HnGMTate@$\TaHm^n(G, M)$}
	Misalkan $G$ grup hingga dan $M$ modul-$G$. Untuk setiap $n \in \Z$,
	definisikan
% segment-id: o014.li2.chapter6.tate-cohomology.g001
	\begin{align*}
		\TaHm^n(G, M) & := \Hm^n(G, M) \quad (n \geq 1), \\
		\TaHm^0(G, M) & := M^G / \nu M \\
		& = \Coker\left[ \nu: M_G \to M^G \right] \twoheadleftarrow \Hm^0(G, M), \\
		\TaHm^{-1}(G, M) & := M^{\nu=0} / \mathfrak{I} M \\
		& = \Ker\left[ \nu: M_G \to M^G \right] \hookrightarrow \Hm_0(G, M), \\
		\TaHm^{-n}(G, M) & := \Hm_{n-1}(G, M) \quad (n \geq 2).
	\end{align*}
	Semua modul ini funktorial terhadap $M$.
\end{definition}

% segment-id: o014.li2.chapter6.tate-cohomology.p004
Kohomologi Tate sering digunakan dalam teori Galois dan teori bilangan
aljabar; salah satu contohnya ialah $\TaHm^{-1}$ yang muncul dalam
\cite[Definisi 9.6.3]{Li1}. Gagasan ini berasal dari J.\ Tate.

% segment-id: o014.li2.chapter6.tate-cohomology.p005
Teori kohomologi yang berguna biasanya memiliki barisan eksak panjang, dan
kohomologi Tate pun demikian.

% segment-id: o014.li2.chapter6.tate-cohomology.t001
\begin{theorem}\label{prop:Tate-coh-long}
	Misalkan $0 \to M' \to M \to M'' \to 0$ barisan eksak pendek modul-$G$.
	Maka terdapat barisan eksak panjang
% segment-id: o014.li2.chapter6.tate-cohomology.q001
	\[ \cdots \to \TaHm^n(G, M) \to \TaHm^n(G, M'') \xrightarrow{\delta^n} \TaHm^{n+1}(G, M') \to \TaHm^{n+1}(G, M) \to \cdots \]
	yang membentang tak terhingga ke kedua arah. Homomorfisme penghubung
	$\delta^n$ bersifat funktorial terhadap morfisme antarbarisan eksak
	pendek, yaitu terhadap diagram komutatif dengan baris-baris eksak
% segment-id: o014.li2.chapter6.tate-cohomology.g002
	\[\begin{tikzcd}
		0 \arrow[r] & M' \arrow[d] \arrow[r] & M \arrow[d] \arrow[r] & M'' \arrow[d] \arrow[r] & 0 \\
		0 \arrow[r] & L' \arrow[r] & L \arrow[r] & L'' \arrow[r] & 0
	\end{tikzcd}\]
\end{theorem}

% segment-id: o014.li2.chapter6.tate-cohomology.pf001
\begin{proof}
	Sambungkan barisan eksak panjang homologi dan kohomologi sebagai berikut:
% segment-id: o014.li2.chapter6.tate-cohomology.g003
	\[\begin{tikzcd}[column sep=small, every cell/.append style={font = \footnotesize}]
		\cdots \arrow[r] & \Hm_1(G, M'') \arrow[r] \arrow[d] & \Hm_0(G, M') \arrow[r] \arrow[d, "{N_1}"] & \Hm_0(G, M) \arrow[r] \arrow[d, "N_2"] & \Hm_0(G, M'') \arrow[r] \arrow[d, "N_3"] & 0 \arrow[d] & \\
		& 0 \arrow[r] & \Hm^0(G, M') \arrow[r] & \Hm^0(G, M) \arrow[r] & \Hm^0(G, M'') \arrow[r] & \Hm^1(G, M') \arrow[r] & \cdots
	\end{tikzcd}\]
	Di sini, $N_1,N_2,N_3$ semuanya merupakan homomorfisme yang sebelumnya
	ditulis $\nu$. Diagram di atas komutatif dan baris-barisnya eksak; satu-satunya
	hal yang perlu dijelaskan ialah komutativitas kedua persegi yang menyentuh
	$0$. Hal ini dapat diperiksa secara konkret dengan menggunakan kompleks
	standar dan kompleks rantai, serta mengingat kembali morfisme penghubung
	$\Hm_1 \to \Hm_0$ dan $\Hm^0 \to \Hm^1$; verifikasi ini diserahkan sebagai
	latihan.

	Lema Ular kemudian memberikan barisan eksak
% segment-id: o014.li2.chapter6.tate-cohomology.g004
	\[\begin{tikzcd}[row sep=small, column sep=small, every cell/.append style={font = \footnotesize}]
		\Ker(N_1) \arrow[r] & \Ker(N_2) \arrow[r] & \Ker(N_3) \arrow[r, "{\delta^{-1}}"] & \Coker(N_1) \arrow[r] & \Coker(N_2) \arrow[r] & \Coker(N_3) \\
		\TaHm^{-1}(G, M') \arrow[equal, u] & \TaHm^{-1}(G, M) \arrow[equal, u] & \TaHm^{-1}(G, M'') \arrow[equal, u] & \TaHm^0(G, M') \arrow[equal, u] & \TaHm^0(G, M) \arrow[equal, u] & \TaHm^0(G, M'') \arrow[equal, u]
	\end{tikzcd}\]
	Inilah satu bagian dari barisan eksak yang dicari, sedangkan funktorialitas
	homomorfisme penghubung $\delta^{-1}$ sudah dijamin oleh Lema Ular.
	Bagian-bagian lainnya tereduksi ke barisan eksak panjang homologi grup dan
	kohomologi grup, serta diagram komutatif dengan baris-baris eksak yang
	dibuktikan pada langkah pertama.
\end{proof}

% segment-id: o014.li2.chapter6.tate-cohomology.p006
Selanjutnya kita nyatakan versi Lema Shapiro (Teorema~\ref{prop:Shapiro}) untuk
kohomologi Tate. Ingat bahwa jika $G$ hingga, maka bagi setiap subgrup $H$ kedua
funktor induksi
$\Ind^G_H,\iInd^G_H:H\dcate{Mod}\to G\dcate{Mod}$ berimpit; lihat
Korolari~\ref{prop:iInd-Ind}. Mulai sekarang, keduanya ditulis $\Ind^G_H$.

% segment-id: o014.li2.chapter6.tate-cohomology.pr001
\begin{proposition}
% segment-id: o014.li2.chapter6.tate-cohomology.x003
	\index{Lema Shapiro}
	Untuk grup hingga $G$, subgrup $H$, dan modul-$H$ $N$, terdapat
	isomorfisme kanonik
% segment-id: o014.li2.chapter6.tate-cohomology.q002
	\[ \TaHm^n(H, N) \simeq \TaHm^n\left( G, \Ind^G_H(N)\right), \quad n \in \Z. \]
\end{proposition}

% segment-id: o014.li2.chapter6.tate-cohomology.pf002
\begin{proof}
	Untuk $n \geq 1$ atau $n \leq -2$, pernyataan ini tereduksi ke
	Teorema~\ref{prop:Shapiro}. Selanjutnya, tangani kasus $n=-1,0$.
	Realisasikan $\Ind^G_H(N)$ sebagai ruang pemetaan melalui
	Lema~\ref{prop:Ind-as-mappings}. Tuliskan $\nu^G$ (atau $\nu^H$) untuk
	$\nu$ yang bersesuaian dengan $G$ (atau $H$). Masalahnya tereduksi ke
	pembuktian bahwa diagram berikut komutatif:
% segment-id: o014.li2.chapter6.tate-cohomology.g005
	\[\begin{tikzcd}
		\Ind^G_H(N)_G \arrow[r, "{\nu^G}"] & \Ind^G_H(N)^G \arrow[d, "\sim" sloped] \\
		N_H \arrow[u, "\sim" sloped] \arrow[r, "{\nu^H}"'] & N^H
	\end{tikzcd}\]
	Isomorfisme-isomorfisme vertikal berasal dari
	Korolari~\ref{prop:Ind-invariant}. Misalkan $y\in N$, dan tuliskan $[y]$
	untuk citranya dalam $N_H$. Citra $[y]$ dalam $\Ind^G_H(N)_G$ ialah $[f]$,
	dengan $f\in\Ind^G_H(N)$ memenuhi $f(1_G)=y$ dan
	$\Supp(f)\subset H$. Dengan demikian, citra
	$\nu^Gf:x\mapsto\sum_{g\in G}f(xg)$ dalam $N^H$ ialah
	$(\nu^Gf)(1_G)=\sum_gf(g)$. Akan tetapi, ini juga sama dengan
	$\sum_{h\in H}f(h)=\sum_hhf(1_G)=\nu^Hy$. Selesai.
\end{proof}

% segment-id: o014.li2.chapter6.tate-cohomology.p007
Argumen tersebut dapat diperkuat untuk menunjukkan bahwa isomorfisme kanonik
ini kompatibel dengan barisan eksak panjang pada kedua ruas.

% segment-id: o014.li2.chapter6.tate-cohomology.c001
\begin{corollary}\label{prop:Tate-ind-vanishing}
	Misalkan $G$ grup hingga. Untuk setiap modul-$\Bbbk$ $N$, berlaku
% segment-id: o014.li2.chapter6.tate-cohomology.q003
	\[ \TaHm^n\left(G, \Ind^G_{\{1\}}(N)\right) = 0, \quad n \in \Z. \]
\end{corollary}

% segment-id: o014.li2.chapter6.tate-cohomology.pf003
\begin{proof}
	Jelas bahwa $\TaHm^n(\{1\},N)=0$ untuk setiap $n\in\Z$.
\end{proof}

% segment-id: o014.li2.chapter6.tate-cohomology.ex001
\begin{example}\label{eg:additive-90-finite}
% segment-id: o014.li2.chapter6.tate-cohomology.x004
	\index{Teorema Hilbert 90}
	Misalkan $E|F$ perluasan Galois hingga dari medan-medan, dengan grup Galois
	$\Gal(E|F)$ beraksi pada $E$. Ambil gelanggang koefisien
	$\Bbbk=F$ atau $\Z$. Teorema Basis Normal
	\cite[Teorema 9.5.6]{Li1} setara dengan pernyataan
	$E\simeq\Ind^{\Gal(E|F)}_{\{1\}}(F)$, sehingga
	Korolari~\ref{prop:Tate-ind-vanishing} mengakibatkan
% segment-id: o014.li2.chapter6.tate-cohomology.q004
	\[ \TaHm^n(\Gal(E|F), E) = 0, \quad n \in \Z; \]
	oleh sebab itu,
	$n\geq1\implies\Hm^n(\Gal(E|F),E)=0$. Jika $E|F$ merupakan perluasan
	siklik dan $n=-1$, hal ini memberikan bukti lain untuk versi aditif
	Teorema Hilbert 90 yang disebutkan dalam
	\cite[Teorema 9.6.10]{Li1}. Untuk versi multiplikatif, lihat
	Contoh~\ref{eg:multiplicative-90-finite}.
\end{example}

% segment-id: o014.li2.chapter6.tate-cohomology.p008
Untuk $H\subset G$ umum dan modul-$G$ $M$, adjoin kiri dan kanan dari restriksi
memberikan morfisme kanonik
% segment-id: o014.li2.chapter6.tate-cohomology.q005
\[ \epsilon_M: \iInd^G_H(\Res^G_H(M)) \to M, \quad \eta_M: M \to \Ind^G_H(\Res^G_H(M)). \]

% segment-id: o014.li2.chapter6.tate-cohomology.p009
Berdasarkan deskripsi konkret pasangan adjoin
(Proposisi~\ref{prop:pullback-two-adjoint}),
$\epsilon_M(g\otimes x)=gx$, sedangkan $\eta_M(x)=[g\mapsto gx]$.
Karena itu, $\epsilon_M$ surjektif dan $\eta_M$ injektif. Kembali ke kasus grup
hingga, diperoleh
$\iInd^G_{\{1\}}=\Ind^G_{\{1\}}$. Dengan menggunakan $\epsilon_M$ dan
$\eta_M$, sifat lenyap dalam Korolari~\ref{prop:Tate-ind-vanishing}
menunjukkan bahwa setiap $\TaHm^n(G,\cdot)$ dapat dihapus sekaligus dapat
dihapus secara dual (Definisi~\ref{def:erasable}). Oleh karena itu, teknik
pergeseran dimensi dalam Proposisi~\ref{prop:dimension-shifting} mempunyai
rumusan sederhana berikut.

% segment-id: o014.li2.chapter6.tate-cohomology.c002
\begin{corollary}[Pergeseran dimensi]\label{prop:Tate-dim-shift}
	Untuk setiap modul-$G$ $M$, seperti biasa tuliskan
	$M$ sebagai singkatan dari $\Res^G_{\{1\}}M$, dan tinjau barisan-barisan
	eksak pendek
% segment-id: o014.li2.chapter6.tate-cohomology.g006
	\begin{gather*}
		0 \to M' \to \Ind^G_{\{1\}}(M) \xrightarrow{\epsilon_M} M \to 0, \\
		0 \to M \xrightarrow{\eta_M} \Ind^G_{\{1\}}(M) \to M'' \to 0.
	\end{gather*}
	Homomorfisme penghubung yang bersesuaian dalam barisan eksak panjang
	memberikan isomorfisme kanonik
% segment-id: o014.li2.chapter6.tate-cohomology.q006
	\[ \TaHm^{n-1}(G, M'') \rightiso \TaHm^n(G, M) \rightiso \TaHm^{n+1}(G, M'), \quad n \in \Z. \]
\end{corollary}

% segment-id: o014.li2.chapter6.tate-cohomology.pf004
\begin{proof}
	Gabungkan barisan eksak panjang
	(Teorema~\ref{prop:Tate-coh-long}) dengan sifat lenyap dalam
	Korolari~\ref{prop:Tate-ind-vanishing}.
\end{proof}

% segment-id: o014.li2.chapter6.tate-cohomology.p010
Banyak sifat kohomologi Tate dapat direduksi melalui pergeseran dimensi ke
kasus $n=0$.

% segment-id: o014.li2.chapter6.tate-cohomology.p011
Kembali ke definisi. Pilih sembarang resolusi projektif $P\to M$ dan resolusi
injektif $M\to I$ dari modul-$G$ $M$, keduanya dipandang sebagai kompleks, lalu
tinjau morfisme kompleks $P_G\to I^G$, yang dapat ditampilkan sebagai diagram
% segment-id: o014.li2.chapter6.tate-cohomology.g007
\[\begin{tikzcd}
	\cdots \arrow[r] & 0 \arrow[r] & 0 \arrow[r] & I^{0, G} \arrow[r] & I^{1, G} \arrow[r] & I^{2, G} \arrow[r] & \cdots \\
	\cdots \arrow[r] & P_{2, G} \arrow[r] \arrow[u] & P_{1, G} \arrow[r] \arrow[u] & P_{0, G} \arrow[u, "\Omega"] \arrow[r] & 0 \arrow[r] \arrow[u] & 0 \arrow[r] \arrow[u] & \cdots
\end{tikzcd}\]
Di sini, $\Omega$ merupakan komposisi
$P_{0,G}\twoheadrightarrow M_G\xrightarrow{\nu}M^G\hookrightarrow I^{0,G}$;
komutativitasnya jelas. Ambil kerucut pemetaan
$\Cone\left[P_G\to I^G\right]$, atau dengan kata lain kompleks total dari
diagram di atas (dengan $P_G$ ditempatkan pada koordinat vertikal $-1$).
Pembaca diminta memverifikasi bahwa
% segment-id: o014.li2.chapter6.tate-cohomology.e002
\begin{equation}\label{eqn:Tate-cone-1}
	\Hm^n\left( \Cone\left[P_G \to I^G\right] \right) = \TaHm^n(G, M).
\end{equation}

% segment-id: o014.li2.chapter6.tate-cohomology.p012
Persamaan di atas sebenarnya menunjukkan cara mendefinisikan morfisme kanonik
$\Bbbk \otimesL[{\Bbbk[G]}] M \to \RHom_G(\Bbbk, M)$ dalam kategori turunan
$\cate{D}(G)$ dari modul-$G$. Lengkapi morfisme ini menjadi segitiga
terbedakan di $\cate{D}(G)$,
% segment-id: o014.li2.chapter6.tate-cohomology.q007
\[ \Bbbk \otimesL[{\Bbbk[G]}] M \to \RHom_G(\Bbbk, M) \to \mathrm{Tate}(G, M) \xrightarrow{+1}, \]
sehingga
$\Hm^n(\mathrm{Tate}(G,M))\simeq\TaHm^n(G,M)$. Masalah konstruksi ini ialah
bahwa pilihan-pilihan yang berbeda untuk $\mathrm{Tate}(G,M)$ dalam segitiga
terbedakan saling isomorfik, tetapi isomorfisme-isomorfisme tersebut tidak
kanonik. Ini merupakan salah satu kekurangan utama teori kategori turunan;
lihat Korolari~\ref{prop:Cone-uniqueness} beserta penjelasan sesudahnya.

% segment-id: o014.li2.chapter6.tate-cohomology.p013
Namun, andaikan
$\Bbbk \otimesL[{\Bbbk[G]}] M$ (atau $\RHom_G(\Bbbk,M)$) direpresentasikan
oleh suatu kompleks terpilih yang terkonsentrasi pada derajat takpositif (atau
taknegatif), termasuk tetapi tidak terbatas pada $P_G$ (atau $I^G$) dalam
\eqref{eqn:Tate-cone-1}. Kita tetap dapat mendefinisikan morfisme $\Omega$ dan
mengambil kerucut pemetaan sebagai $\mathrm{Tate}(G,M)$, lalu menafsirkannya
sebagai kokernel homotopi dalam arti
Catatan~\ref{rem:homotopy-kernel-cokernel}.\footnote{Mengambil kernel homotopi
dari $\Bbbk \otimesL[{\Bbbk[G]}] M \to \RHom_G(\Bbbk, M)$ tidak menghasilkan
struktur baru; kernel homotopi itu hanya berbeda dari kokernel homotopi dengan
suatu pergeseran.} Latihan-latihan akan menunjukkan cara menggunakan resolusi
standar dari $\Bbbk$ sebagai gantinya, sehingga menyediakan lebih banyak alat
untuk menghitung kohomologi Tate.

% segment-id: o014.li2.chapter6.tate-cohomology.p014
Sekarang kita bahas lebih lanjut kasus khusus grup siklik hingga.

% segment-id: o014.li2.chapter6.tate-cohomology.ex002
\begin{example}\label{eg:TaHm-periodic}
	Misalkan $m\in\Z_{\geq1}$ dan $C_m$ grup siklik berorde $m$; pilih sembarang
	generator $\sigma$ dari $C_m$. Untuk setiap modul-$C_m$ $A$, kita nyatakan
	bahwa $\TaHm^n(C_m,A)$ ialah $\Hm^n$ dari kompleks berikut
	($n\in\Z$):
% segment-id: o014.li2.chapter6.tate-cohomology.q008
	\[ \cdots \xrightarrow{\sigma-1} A \xrightarrow{\nu} \underbracket{A}_{\text{suku derajat nol}} \xrightarrow{\sigma-1} A \xrightarrow{\nu} A \xrightarrow{\sigma-1} \to \cdots \]
	Kompleks ini berulang dengan periode $2$. Alasannya, bukti
	Proposisi~\ref{prop:group-cohomology-cyclic} telah menuliskan secara konkret
	kompleks-kompleks yang menghitung homologi dan kohomologi; menyambungkan
	keduanya dengan $\nu$ menghasilkan kerucut pemetaan
	$\mathrm{Tate}(G,M)$ yang disebutkan sebelumnya, tepat dalam bentuk di
	atas. Dengan demikian, langsung diperoleh isomorfisme kanonik
% segment-id: o014.li2.chapter6.tate-cohomology.q009
	\[ \TaHm^n(C_m, A) \simeq \begin{cases}
		A^{\nu=0} / (\sigma-1) A, & n\;\text{ganjil} \\
		A^{\sigma=1} / \nu A, & n\;\text{genap}.
	\end{cases}\]

	Periodisitas ini menunjukkan bahwa jika
	$0\to A'\to A\to A''\to0$ merupakan barisan eksak pendek modul-$C_m$,
	maka barisan eksak panjang dalam Teorema~\ref{prop:Tate-coh-long} termuat
	dalam segi enam eksak
% segment-id: o014.li2.chapter6.tate-cohomology.e003
	\begin{equation}\label{eqn:Herbrand-hexagon}
		\begin{tikzcd}[column sep=small]
			& \TaHm^{\text{genap}}(C_m, A') \arrow[r, "f_1"] & \TaHm^{\text{genap}}(C_m, A) \arrow[rd, "f_2"] & \\
			\Hm^{\text{ganjil}}(C_m, A'') \arrow[ru, "f_6"] & & & \TaHm^{\text{genap}}(C_m, A'') \arrow[ld, "f_3"] . \\
			& \TaHm^{\text{ganjil}}(C_m, A) \arrow[lu, "f_5"] & \TaHm^{\text{ganjil}}(C_m, A') \arrow[l, "f_4"] &
		\end{tikzcd}
	\end{equation}
\end{example}

% segment-id: o014.li2.chapter6.tate-cohomology.d002
\begin{definition}
% segment-id: o014.li2.chapter6.tate-cohomology.x005
	\index{hasil bagi Herbrand (Herbrand quotient)}
	Misalkan $A$ modul-$C_m$ ($m\in\Z_{\geq1}$). Dengan asumsi bahwa
	$\TaHm^n(C_m,A)$ hingga untuk setiap $n$, definisikan bilangan rasional
	positif
% segment-id: o014.li2.chapter6.tate-cohomology.q010
	\[ h(C_m, A) := \frac{\left|\TaHm^{\text{genap}}(C_m, A)\right|}{\left|\TaHm^{\text{ganjil}}(C_m, A)\right|}, \]
	yang disebut \emph{hasil bagi Herbrand} dari $A$.
\end{definition}

% segment-id: o014.li2.chapter6.tate-cohomology.p015
Jika $A$ hingga, maka $\TaHm^n(C_m,A)$ jelas hingga. Perhatikan bahwa banyaknya
elemen tidak bergantung pada pilihan gelanggang koefisien $\Bbbk$; kita dapat
memilih $\Bbbk=\Z$. Dalam teori bilangan aljabar, teorema Herbrand berikut
merupakan hasil yang mendasar dan penting.

% segment-id: o014.li2.chapter6.tate-cohomology.t002
\begin{theorem}[J.\ Herbrand]
	Misalkan $m\in\Z_{\geq1}$ dan $C_m$ grup siklik berorde $m$.
% segment-id: o014.li2.chapter6.tate-cohomology.l001
	\begin{enumerate}[(i)]
% segment-id: o014.li2.chapter6.tate-cohomology.i001
		\item Misalkan $0\to A'\to A\to A''\to0$ barisan eksak pendek
		modul-$C_m$. Jika dua dari $h(C_m,A')$, $h(C_m,A)$, dan
		$h(C_m,A'')$ terdefinisi, yang ketiga juga terdefinisi secara otomatis;
		dalam hal ini,
% segment-id: o014.li2.chapter6.tate-cohomology.q011
		\[ h(C_m, A) = h(C_m, A') h(C_m, A''). \]
% segment-id: o014.li2.chapter6.tate-cohomology.i002
		\item Jika $A$ modul-$C_m$ hingga, maka $h(C_m,A)=1$.
	\end{enumerate}
\end{theorem}

% segment-id: o014.li2.chapter6.tate-cohomology.pf005
\begin{proof}
	Untuk~(i), perhatikan segi enam eksak
	\eqref{eqn:Herbrand-hexagon}. Tuliskan
	$n_i:=|\Image(f_i)|$. Dalam arti perkalian kardinalitas, kardinalitas keenam
	suku tersebut berturut-turut adalah
% segment-id: o014.li2.chapter6.tate-cohomology.g008
	\begin{center}\begin{tikzpicture}[commutative diagrams/every diagram,
		one/.style={draw, rectangle, minimum size=6mm, rounded corners=3mm},
		two/.style={draw, rectangle, minimum size=6mm}
		]
		\node[one] (P0) at (60:1.75cm) {$n_1 n_2$};
		\node[two] (P1) at (60+60:1.75cm) {$n_6 n_1$} ;
		\node[one] (P2) at (60+2*60:1.75cm) {$n_5 n_6$};
		\node[two] (P3) at (60+3*60:1.75cm) {$n_4 n_5$};
		\node[one] (P4) at (60+4*60:1.75cm) {$n_3 n_4$};
		\node[two] (P5) at (0:1.75cm) {$n_2 n_3$};
	\end{tikzpicture}\end{center}

	Tiga hasil bagi Herbrand bersesuaian dengan hasil bagi suku-suku
	yang saling berhadapan pada diagram di atas. Jadi, pernyataan~(i) setara
	dengan pengamatan berikut:
% segment-id: o014.li2.chapter6.tate-cohomology.l002
	\begin{itemize}
% segment-id: o014.li2.chapter6.tate-cohomology.i003
		\item jika kardinalitas pada kedua ujung dua diagonal bernilai
		hingga, hal yang sama juga berlaku pada kedua ujung diagonal sisanya;
% segment-id: o014.li2.chapter6.tate-cohomology.i004
		\item hasil kali suku-suku dalam kotak bersudut siku-siku sama dengan
		hasil kali suku-suku dalam kotak bersudut bulat.
	\end{itemize}

	Untuk~(ii), terdapat barisan-barisan eksak pendek dari modul-$\Z$ hingga
% segment-id: o014.li2.chapter6.tate-cohomology.q012
	\[ 0 \to A^{\sigma=1} \to A \xrightarrow{\sigma-1} \mathfrak{I} A \to 0, \quad 0 \to A^{\nu=0} \to A \xrightarrow{\nu} \nu A \to 0; \]
	di sini, $(\sigma-1)A=\mathfrak{I}A$ merupakan penerapan
	Lema~\ref{prop:cyclic-resolution-prep}. Oleh karena itu,
% segment-id: o014.li2.chapter6.tate-cohomology.q013
	\[ |A^{\sigma=1}| \cdot |\mathfrak{I} A| = |A^{\nu=0}| \cdot |\nu A|, \]
	yang dapat disusun ulang menjadi
	$|\TaHm^0(C_m,A)|=|\TaHm^{-1}(C_m,A)|$. Ini membuktikan pernyataan.
\end{proof}
